% Part: sets-functions-relations
% Chapter: relations
% Section: equivalence-relations

\documentclass[../../../include/open-logic-section]{subfiles}

\begin{document}

\olfileid{sfr}{rel}{eqv}

\olsection{د هم ارزښتۍ اړيکې}

په يوه سټ باندې د عينيت اړيکه انعکاسي، متناظره او انتقالي ده۔
هغه اړيکې $R$ چې دا درې واړه خاصيتونه لري، ډېرې عامې دي۔

\begin{defn}[د هم ارزښتۍ اړيکه]
هغه اړيکه $R \subseteq A^2$ چې انعکاسي، متناظره او انتقالي وي،
\emph{د هم ارزښتۍ اړيکه} بللے شي۔ د $A$ !!^{element}s، $x$ او
$y$، د $R$ له مخې \emph{هم ارزښته} بللے شي که $Rxy$ وي۔
\end{defn}

د هم ارزښتۍ اړيکې د \emph{هم ارزښتۍ د ټولګي} مفهوم راپيدا کوي۔
د هم ارزښتۍ يوه اړيکه ساحه په بېلو برخو ``ټوټه کوي''۔ په هره برخه
کښې ټول څيزونه يو له بل سره اړيکه لري؛ او د بېلو برخو څيزونه يو
له بل سره اړيکه نۀ لري۔ کله کله د همدې برخو په اړه \emph{نېغ په نېغه}
خبرې کول ګټور وي۔ د دې دپاره يو تعريف راوړو:

\begin{defn}\ollabel{def:equivalenceclass}
فرض کړئ $R \subseteq A^2$ د هم ارزښتۍ اړيکه ده۔ د هر $x \in A$
دپاره، په $A$ کښې د $x$ \emph{د هم ارزښتۍ ټولګے} دا سټ دے:
$\equivrep{x}{R} = \Setabs{y \in A}{Rxy}$۔ د $R$ له مخې د $A$
\emph{خارج قسمت} دا دے:
$\equivclass{A}{R} = \Setabs{\equivrep{x}{R}}{x \in A}$، يعنې
د هم ارزښتۍ د دې ټولګيو سټ۔
\end{defn}

راتلونکې نتيجه د هم ارزښتۍ د ټولګي د تعريف سموالى ښيي: ثابتوي چې
د هم ارزښتۍ ټولګي په رښتيا د $A$ وېشونکې برخې دي۔

\begin{prop}
که $R \subseteq A^2$ د هم ارزښتۍ اړيکه وي، نو $Rxy$ هله او يوازې
هله چې $\equivrep{x}{R} = \equivrep{y}{R}$ وي۔
\end{prop}

\begin{proof}
د کيڼې خوا نه ښۍ خوا ته د لوري دپاره، فرض کړئ $Rxy$ او
$z \in \equivrep{x}{R}$ وي۔ نو د تعريف له مخې $Rxz$۔ ځکه چې $R$
د هم ارزښتۍ اړيکه ده، $Ryz$۔ (په تفصيل سره: له $Rxy$ او د $R$
له تناظر څخه $Ryx$ لرو؛ او له $Rxz$ او د $R$ له انتقالي والي څخه
$Ryz$ لرو۔) نو $z \in \equivrep{y}{R}$۔ په عمومي توګه،
$\equivrep{x}{R} \subseteq \equivrep{y}{R}$۔ خو بالکل په همدې
ډول، $\equivrep{y}{R} \subseteq \equivrep{x}{R}$۔ نو د غړو له
مخې د برابرۍ د اصل له مخې، $\equivrep{x}{R} = \equivrep{y}{R}$۔

د ښۍ خوا نه کيڼې خوا ته د لوري دپاره، فرض کړئ
$\equivrep{x}{R} = \equivrep{y}{R}$۔ ځکه چې $R$ انعکاسي ده،
$Ryy$، نو $y \in \equivrep{y}{R}$۔ له دې فرضيې څخه چې
$\equivrep{x}{R} = \equivrep{y}{R}$، دا هم لرو چې
$y \in \equivrep{x}{R}$۔ نو $Rxy$۔
\end{proof}

\begin{ex}
د هم ارزښتۍ د اړيکو يوه ښۀ بېلګه د باقي له حسابه راځي۔ د هر $a$،
$b$ او $n \in \PosInt$ دپاره، وايو چې $a \equiv_n b$ هله او
يوازې هله چې د $a$ په $n$ د وېش باقي د $b$ په $n$ د وېش له باقي
سره يو شان وي۔ (په يو څه زياته نښيزه ژبه: $a \equiv_n b$ هله او
يوازې هله چې د کوم $k \in \Int$ دپاره $a - b = kn$ وي۔)
اوس د هر $n$ دپاره $\equiv_n$ د هم ارزښتۍ اړيکه ده۔ او د
$\equiv_n$ له مخې جوړ شوي د هم ارزښتۍ ټولګي په شمېر پوره $n$
بېل ټولګي دي؛ يعنې $\equivclass{\Nat}{\equiv_n}$ په شمېر $n$
!!{element}s لري۔ دا ټولګي دي: د هغو عددونو سټ چې بې له باقي پر
$n$ وېشل کېږي، يعنې $\equivrep{0}{\equiv_n}$؛ د هغو عددونو سټ
چې پر $n$ د وېش باقي ئې $1$ وي، يعنې $\equivrep{1}{\equiv_n}$؛
\ldots؛ او د هغو عددونو سټ چې پر $n$ د وېش باقي ئې $n-1$ وي،
يعنې $\equivrep{n-1}{\equiv_n}$۔
\end{ex}

\begin{prob}
وښيئ چې د هر $n \in \PosInt$ دپاره $\equiv_n$ د هم ارزښتۍ
اړيکه ده، او $\equivclass{\Nat}{\equiv_n}$ په شمېر پوره $n$
غړي لري۔
\end{prob}

\end{document}
